\documentclass[11pt,a4paper]{article}
\usepackage[T1]{fontenc}
\usepackage{lmodern,amsmath,amssymb,booktabs,microtype}
\usepackage[margin=24mm]{geometry}
\usepackage[colorlinks=true,urlcolor=blue,linkcolor=blue]{hyperref}
\newcommand{\EE}{\mathbb E}
\newcommand{\tr}{\operatorname{tr}}
\title{Central charge windows for the memory calculation\\Uniform rank tails and a conserved filter correction}
\author{Alan Day \and Chengheng\thanks{AI research assistant (OpenAI).}}
\date{Local research note --- 6 November 2026}
\begin{document}\maketitle
\begin{abstract}
We extend the positive rank-product sixth-moment estimate from bounded
charge shifts to compact windows of width $\sqrt N$. The leading six
moments are unchanged. In contrast, replacing the sector variance by
its half-filled value inside a growing counterterm leaves a finite
error. The exact two-charge memory filter has a further finite,
charge-dependent $H^2$ correction. These results supply inputs for the
central-window operator calculation; the full memory Gram matrix in
that window is still to be established.
\end{abstract}
\section{The charge window and the literature interface}
Use standard CAR, the full overlapping complex Gaussian $p$-body
ensemble, normalized trace at charge $m$, annealed Hilbert--Schmidt
norm, and $\hbar=1$. The coupling variance $w_N$ is the same in all
charge sectors of a given sample. Put
\begin{equation}
 \lambda_N=p^2/N,\quad m=N/2+xp=N/2+z\sqrt N,
 \quad z=x\sqrt{\lambda_N}.
\end{equation}
All uniform statements have $|x|\le L$ (equivalently bounded $z$),
and $\lambda_N$ in a compact subset of $(0,\infty)$. Integer charge
values are understood. Limits below keep $\lambda_N=\lambda$;
using actual $\lambda_N$ gives the corresponding uniform expansions.

The fixed-charge embedded-ensemble formulation already supplies
the sector variance and reduced spectral parameter~\cite{TT}:
\begin{equation}
 C_m=w_N\binom mp\binom{N-m+p}p,\qquad
 q_m=\exp\!\left[-\frac{p^2N}{m(N-m)}\right].\label{variance}
\end{equation}
The expression for $C_m$ is exact; $q_m$ is the double-scaled
spectral parameter, not an exact finite-$N$ fourth moment.
Fixed-charge and grand-canonical transfer matrices are also treated
in Ref.~\cite{FKN}. We use their separation of charge sectors as
context; no transfer-matrix formula substitutes for a uniform bound
on the marked memory insertions. In particular we keep the common
$w_N$, rather than independently renormalizing the Hamiltonian in
each sector.

\newpage
\section{Uniform sixth moments on a central window}
Let $U,V$ be independent isotropic Hermitian rank-$a,b$ matrices,
normalized by $\EE U^2=\EE V^2=I$, in charge $M$. Set
\begin{equation}
 \pi^{(M)}_{ab}(r)=\EE\tr_M[P_r(UV)]^\dagger P_r(UV),\qquad
 X=x_{a+b}-x_r,\quad x_j=j(N+1-j).
\end{equation}
The existing positive product law is normalized, has $r\le a+b$,
and obeys the Racah product identity~\cite{local}:
\begin{align}
 \sum_r\pi^{(M)}_{ab}(r)\chi_r(n)&=\chi_a(n)\chi_b(n),\\
 \chi_a(n)&={}_4F_3\!\left[\begin{matrix}-a,a-N-1,-n,n-N-1\\
                 -M,M-N,-N-1\end{matrix};1\right].
\end{align}
Write $P_j(t)=\prod_{\ell<j}(t-x_\ell)$ and
\begin{equation}
 c_{nj}=\frac{(-n)_j(n-N-1)_j}
 {(-M)_j(M-N)_j(-N-1)_j j!}.
\end{equation}
Then $F_j=\EE_\pi P_j(x_r)$ satisfies the exact triangular recursion
\begin{equation}
 F_0=1,\qquad
 F_n=\frac{\chi_a(n)\chi_b(n)-\sum_{j<n}c_{nj}F_j}{c_{nn}}.
 \label{recurrence}
\end{equation}
Set $\epsilon=N^{-1/2}$, $a=A/\epsilon$, $b=B/\epsilon$,
and $M=1/(2\epsilon^2)+z/\epsilon$. In the scaled variable
$Y=\epsilon^3x_r$, the coefficient $c_{nj}\epsilon^{-3j}$ is
\begin{equation}
 (-4)^j\binom nj\epsilon^j
 \prod_{k=0}^{j-1}
 \frac{1-(n-1+k)\epsilon^2}
 {(1+2z\epsilon-2k\epsilon^2)
  (1-2z\epsilon-2k\epsilon^2)(1+(1-k)\epsilon^2)}.
 \label{centralcoef}
\end{equation}
This is the substantive change from a bounded shift $M-N/2$.
The roots are $\epsilon^3x_\ell=\ell\epsilon+
\ell(1-\ell)\epsilon^3$ and the input is
$Y_a=A-A^2\epsilon+A\epsilon^2$.

The exact series recursion over $\mathbb Q[A,B,z]$ through
$\epsilon^{12}$ cancels all 21 negative Laurent coefficients in
the first six moments of $X/N$. Its constant terms give
\begin{equation}
 \EE_\pi(X/N)^j=T_j(2AB)+O(\epsilon),\quad 1\le j\le6,
 \qquad T_{j+1}(t)=t[T_j(t)+T'_j(t)],\ T_0=1.\label{moments}
\end{equation}
After extracting the known powers of $\epsilon$, each divisor in
\eqref{recurrence} has nonzero constant term $(-4)^n$.
All other denominators in \eqref{centralcoef} remain nonzero on
compact $A,B,z$ sets for sufficiently small $\epsilon$.
Consequently the exact rational functions are analytic there and
the remainders in \eqref{moments} are uniform, not merely formal.

\newpage
\section{What the tail bound controls}
The sixth polynomial is
\begin{equation}
 T_6(t)=t^6+15t^5+65t^4+90t^3+31t^2+t.
\end{equation}
For $a,b\le p$, the output deficit $S=a+b-r$ is nonnegative and
\begin{equation}
 X=S[N+1-2(a+b)+S]\ge S(N+1-4p).
\end{equation}
Thus, uniformly on the central window and on its adjacent charges,
\begin{equation}
 \sup_{a,b\le p,\ |M-N/2|\le L\sqrt N}\EE_\pi S^6
 \le K_{L,\lambda},\qquad
 \sum_{S>R}\pi_{ab}^{(M)}(r)\le K_{L,\lambda}R^{-6}.
 \label{tail}
\end{equation}
The norm of an independent rank-product tail is the square root
of this probability. The constant is allowed to depend on $L$;
the calculation does not cover a growing $L=L_N$.

This distribution refers to independent isotropic rank inputs.
The actual memory operator contains coherent products in adjacent
charges and ordinary Gaussian chaos subtractions. Applying
\eqref{tail} to that operator still requires the appropriate
charged norm-transfer and input-rank estimates. It is not a direct
bound on the whole memory remainder.

\paragraph{Executed checks.}
The polynomial-ring calculation verifies six constant terms and
21 cancellations. The four full finite distributions at
\[
 (N,M,a,b)=(16,6,2,3),(16,10,2,3),(20,7,3,3),(20,13,3,3)
\]
independently reproduce all 28 moments from order zero to six.
Their $S^6$ bounds are checked as exact fractions. Forty further
exact-rational cases use $\lambda=1/4$, $p=16,32,64,128$,
$z=-2,-1,0,1,2$, and input deficits $(0,0),(1,2)$; reflection
$M\leftrightarrow N-M$ is checked explicitly.

For example, at $p=128$, $z=2$, $a=b=p$,
$\EE(X/N)^6=22.4617786805$, while $T_6(1/2)=24.046875$.
The finite difference is retained; these are convergence checks,
not equality to the limiting moments at finite $N$.

\newpage
\section{A vanishing variance difference leaves a finite counterterm}
Let $C_0=C_{N/2}$ for even $N$. An exact product from
\eqref{variance} is
\begin{equation}
 \frac{C_m}{C_0}=
 \prod_{i=0}^{p-1}
 \left(1+\frac{xp}{N/2-i}\right)
 \left(1-\frac{xp}{N/2+p-i}\right).
\end{equation}
With $\eta=1/p$ and $y=i/p$, the logarithm of each paired
factor has the uniform expansion
\begin{align}
 &\log\left(1+\frac{2\lambda x\eta}{1-2\lambda y\eta}\right)
 +\log\left(1-\frac{2\lambda x\eta}{1+2\lambda(1-y)\eta}\right)
 \notag\\
 &\hspace{25mm}=4\lambda^2x(1-x)\eta^2+O(\eta^3).
\end{align}
Summing $p$ terms and exponentiating gives
\begin{equation}
 \boxed{\quad p\left(\frac{C_m}{C_0}-1\right)
 =4\lambda^2x(1-x)+O(p^{-1}).\quad}\label{drift}
\end{equation}
The asymmetric center is not a mistake: the exact variance obeys
$C_m=C_{N+p-m}$, corresponding to $x\mapsto1-x$.
It is maximal near $m=(N+p)/2$, not exactly at half filling.
The original ensemble contains all overlapping index sets; no
particle-hole normal-ordering modification has been made.

Meanwhile
\begin{equation}
 \log q_m=-\frac{4\lambda}{1-4\lambda^2x^2/p^2}
 =-4\lambda-\frac{16\lambda^3x^2}{p^2}+O(p^{-4}).
\end{equation}
Thus a common leading spectral shape can coexist with a finite
correction to a subtraction proportional to $pC_m$.
Replacing $C_m$ by $C_0$ in that subtraction is not justified by
$C_m/C_0\to1$.

Thirty exact binomial ratios, $p=16,32,64,128,256$ and
$x=-2,-1,0,1/2,1,2$, verify the reflection identity and record
the deviation from \eqref{drift}. These use a common coupling
variance; no independent sector rescaling is performed.

\newpage
\section{The raw core filter also acquires a finite charge correction}
Use the exact previous two-charge decomposition~\cite{local},
with $A_{ab;r}=Z^-_{ab;r}+(2p-m)Z^+_{ab;r}$:
\begin{equation}
 Q_{\mathcal E}=\sum_{a\le b,r}
 (g_H Z^+_{ab;r}+g_A A_{ab;r}).
\end{equation}
Here $Q_{\mathcal E}$ is the quadratic contraction of the raw
quartic core, prior to its scalar centering. Its known rational
coefficients can be organized as
\begin{equation}
 g_H/C_m=B_0+B_a\kappa_a/C_m+B_b\kappa_b/C_m+B_r\kappa_r/C_m,
 \quad \kappa_j=w_N\binom{m-j}p\binom{N-m+p-j}p.
\end{equation}
For $a=p-u$, $b=p-v$, $r=2p-u-v-S$, fixed deficits, and
$m=N/2+xp$, the exact rational expansions are
\begin{align}
 B_0&=2p+\lambda+4\lambda x+O(p^{-1}),\\
 B_a=B_b&=-2p-5\lambda-4\lambda x+O(p^{-1}),\\
 B_r&=2p+9\lambda+4\lambda x+O(p^{-1}).
\end{align}
For $j=tp-s$, the paired binomial product gives
\begin{equation}
 \log(\kappa_j/C_m)=-4\lambda t+
 (4\lambda s-4\lambda^2t^2)/p+O(p^{-2}).
\end{equation}
The charge coordinate cancels through this order. Therefore,
with $q=e^{-4\lambda}$,
\begin{align}
 g_H/C_m={}&2p(1-q)^2+K_0+4\lambda x(1-q)^2\notag\\
 &+8\lambda(q^2-q)(u+v)+8\lambda q^2S+O(p^{-1}),\label{gH}\\
 K_0={}&\lambda[1-10q+9q^2+16\lambda q-32\lambda q^2].
\end{align}
The additional term is rank-independent. After ordinary Gaussian
centering it multiplies $C_m(h_m^2-C_mI)$, a conserved direction.
It must be included before studying the remaining dynamic field.

For $g_A/C_m$ the four rational prefactors, at arbitrary bounded
rank proportions, remain
\begin{equation}
 (3/2,-1,-1,1/2)+\frac{\lambda}{p}(8,-5/2,-5/2,3)+O(p^{-2}).
\end{equation}
In particular the fixed-deficit leading coefficient is still
$(3-4q+q^2)/2$. Uniform operator convergence additionally needs
the coherent norm estimates identified on page 3.

\newpage
\section{Verification and the next operator calculation}
The saved finite five-word implementation is independent of the
rational-coefficient expansion. At $\lambda=1/4$, five values
$p=16,32,64,128,256$, $x=-1,0,1$, and two deficit triples,
it verifies 60 exact rational identities for $g_H$ and $g_A$.
The normalized difference
\begin{equation}
 g_H(N/2+xp)/C_m-g_H(N/2)/C_0
 \longrightarrow4\lambda x(1-q)^2
\end{equation}
is evaluated at 80 and 120 decimal digits with agreement better
than $10^{-70}$ in the stored quantities. For $p=256$, $x=1$,
$(u,v,S)=(2,1,3)$, its value is $0.39841612746439$;
the limiting value is $(1-e^{-1})^2=0.39957640089373$.
The finite discrepancy is recorded, not removed by fitting.
Zero-limit cases use absolute errors. These are exact-rational
inputs followed by high-precision evaluations, not interval arithmetic.

\paragraph{Reproducibility.}
The three records are \texttt{central-moments.json},
\texttt{central-scale.json}, and \texttt{central-filter-checks.json}
under \texttt{results/susskind-memory-central-window}.
They retain parameter choices, dependency versions, and input hashes.
The first symbolic run exposed a reused dummy variable; the first
scale check used structural rather than algebraic symbolic equality.
Both initial sources and explanations are preserved. Corrections
kept exact algebraic assertions. Previous scientific records were
not overwritten.

\paragraph{Continuation.}
The central-window tail estimate and the finite conserved
correction are now explicit. Next derive the actual charged
norm transfer and marked overlap estimates uniformly on this
window, including the $A$ input-rank tail, then assemble the
memory Gram matrix with sector-dependent conserved subtractions.
Compact-window control by itself does not justify summing the
whole charge distribution; an integrable bound outside the window
is also required. After that, control the two-insertion finite-memory
remainder and connect the microscopic fields to the Susskind
singlet/meson spatial and mass dictionary.

\begin{thebibliography}{9}\small
\bibitem{TT} J. Tall and S. Tomsovic,
\href{https://arxiv.org/html/2604.14522v2}{arXiv:2604.14522v2},
Secs.~2.1 and 6: fixed-sector normalization and reduced moments.
\bibitem{FKN} S. F\"orste, Y. Kruse and S. Natu,
\href{https://arxiv.org/html/2512.07715v2}{arXiv:2512.07715v2},
Secs.~3.1--3.2: canonical and grand-canonical transfer matrices.
\bibitem{local} A. Day and Chengheng, local notes:
\emph{An exact two-charge formula for the memory core};
\emph{The quadratic memory core in the double-scaled limit};
\emph{Restoring the dynamic memory descendants} (2026-11-06).
\end{thebibliography}
\end{document}
